\section{Domain Invariance Experiment}
\label{sec:domain}
% ====================================================================

The preceding results establish that DA2-based depth estimation
enables strong parkour performance.  We now test the core hypothesis:
does this pipeline provide domain invariance?

\subsection{Hypothesis}

\begin{itemize}
  \item $H_1$: The DA2-based student maintains performance under
    visual domain shifts (texture, lighting, color), because the depth
    estimator maps diverse appearances to a canonical depth space.
  \item $H_0$: The DA2 student degrades significantly under
    visual perturbations.
\end{itemize}

\subsection{Visual Perturbations}

We apply perturbations at \emph{test time only}---the student was
trained exclusively with checkerboard-textured terrain.  Perturbations
span two categories:

\paragraph{Texture replacement.}
Replace the training checkerboard with six alternatives: no texture
(untextured geometry), solid colors (red, blue), flat grey, perlin
noise, or a brick pattern.  These range from completely featureless
surfaces (no texture, solid colors, flat grey) to structured
alternatives (noise, bricks).

\paragraph{Color jitter.}
Apply image-space perturbations to the rendered RGB frame before DA2
inference: global brightness scaling ($\pm30\%$, $\pm60\%$) and
per-channel color shifts ($\pm15\%$, $\pm30\%$).  These simulate
lighting and white-balance variation without changing geometry.

\paragraph{Combined.}
Simultaneous texture replacement + color jitter at two severity
levels: mild (bricks + brightness 30\% + color 15\%) and extreme
(noise + brightness 60\% + color 30\%).

\subsection{Protocol}

The DA2-base+texture student (Section~\ref{sec:da2_results}) is
evaluated under all 13 conditions with 128 episodes on forced gap
terrain (Row~0).  We report success rate, progress ratio, and
the delta from baseline (training visual conditions).

\subsection{Results}

Table~\ref{tab:domain_inv} reports the full results.  The
perturbations reveal a clean separation between two types of
visual domain shift.

\begin{table}[h]
  \centering
  \caption{Domain invariance results: DA2-base+texture student on gap
  Row~0 (128 episodes per condition).}
  \label{tab:domain_inv}
  \small
  \begin{tabular}{@{}p{3.6cm}rrr@{}}
    \toprule
    Perturbation & Succ. & Prog. & $\Delta$ \\
    \midrule
    Baseline (checkerboard) & 96.9\% & .948 & --- \\
    \midrule
    \multicolumn{4}{@{}l}{\textit{Texture replacement (structured)}} \\
    \quad Bricks & 89.8\% & .908 & $-$7.1 \\
    \quad Noise & 78.1\% & .843 & $-$18.8 \\
    \midrule
    \multicolumn{4}{@{}l}{\textit{Texture removal (featureless)}} \\
    \quad No texture & 18.0\% & .437 & $-$78.9 \\
    \quad Solid red & 23.4\% & .471 & $-$73.5 \\
    \quad Solid blue & 14.8\% & .417 & $-$82.1 \\
    \quad Flat grey & 18.0\% & .456 & $-$78.9 \\
    \midrule
    \multicolumn{4}{@{}l}{\textit{Color / brightness jitter}} \\
    \quad Brightness $\pm$30\% & 95.3\% & .937 & $-$1.6 \\
    \quad Brightness $\pm$60\% & 94.5\% & .935 & $-$2.4 \\
    \quad Color $\pm$15\% & 91.4\% & .914 & $-$5.5 \\
    \quad Color $\pm$30\% & 76.6\% & .816 & $-$20.3 \\
    \midrule
    \multicolumn{4}{@{}l}{\textit{Combined}} \\
    \quad Mild (bricks+b30+c15) & 82.0\% & .862 & $-$14.9 \\
    \quad Extreme (noise+b60+c30) & 63.3\% & .747 & $-$33.6 \\
    \bottomrule
  \end{tabular}
\end{table}

\subsection{Analysis}

The results reveal that the DA2 pipeline provides \emph{partial}
domain invariance, with a clear dichotomy between two classes of
visual perturbation.

\paragraph{Color and brightness changes.}
DA2 is highly robust to image-space color and brightness perturbations
---the kind of variation most relevant to sim-to-real transfer.
Brightness jitter at $\pm$60\% causes only a 2.4~pp drop;
color jitter at $\pm$15\% causes 5.5~pp.  These shifts simulate
the lighting and white-balance differences between a simulated
scene and a real environment.  Under these perturbations, the DA2
student clearly satisfies the domain-invariance criterion ($<$10~pp
drop) at moderate severity.

\paragraph{Texture replacement.}
Replacing the training texture with a different \emph{structured}
texture (bricks: $-$7.1~pp; noise: $-$18.8~pp) produces moderate
degradation.  The bricks result is within the 10~pp invariance
threshold, suggesting DA2 generalizes across structured visual
patterns.  However, removing texture entirely---leaving featureless
surfaces (no texture, solid colors, flat grey)---causes catastrophic
failure ($-$73 to $-$82~pp).

\paragraph{Interpreting the texture sensitivity.}
The texture-removal failure is \emph{not} a domain invariance
failure of DA2; rather, it is an input quality requirement.
Monocular depth estimation fundamentally relies on visual features
(texture gradients, edges, perspective cues) to infer geometry from
a single image.  A uniformly colored, featureless surface provides
no visual information from which to estimate depth---just as a
depth camera would fail on a perfectly transparent or specular
surface.  Critically, real-world surfaces (concrete, wood, soil,
asphalt) universally have natural texture, so this failure mode
is an artifact of the simulation's ability to create unrealistically
featureless geometry.

\paragraph{Combined perturbations.}
The combined-mild condition (bricks + brightness 30\% + color 15\%)
maintains 82.0\% success ($-$14.9~pp), while the extreme condition
(noise + brightness 60\% + color 30\%) drops to 63.3\% ($-$33.6~pp).
This shows that perturbation effects compound, but the DA2 student
degrades gracefully rather than catastrophically even under
simultaneous texture replacement and aggressive color jitter.
